\documentclass[11pt]{article}
\usepackage{amsmath,amssymb,amsthm}
\usepackage[margin=1in]{geometry}

\newtheorem{theorem}{Theorem}
\newtheorem{lemma}[theorem]{Lemma}
\newtheorem{definition}[theorem]{Definition}
\newtheorem{remark}[theorem]{Remark}

\title{Formalization of G\"odel's Second Incompleteness Theorem\\in Guard-style Recursive Arithmetic}
\author{Summary of Agda development}
\date{April 2026}

\begin{document}
\maketitle

\section{Overview}

We formalize G\"odel's second incompleteness theorem in the style of
Rose~\cite{rose1961}, using a system of binary trees with 1-ary and 2-ary
functors (Fun1, Fun2), an equational derivation system (Deriv), and
tree recursion (Rec). The formalization is in Agda with
\texttt{--safe --without-K --exact-split}. No postulates are used.

The main result:

\begin{theorem}[godelII]
For any hypothesis \texttt{hyp}:
\[
  \mathit{ProofE}\;\mathit{hyp} \to \mathit{Consistent}\;\mathit{hyp} \to
  \mathit{Deriv}\;\mathit{hyp}\;\mathit{godelSentence} \to \mathit{Empty}
\]
\end{theorem}

\noindent where:
\begin{itemize}
\item $\mathit{ProofE}\;\mathit{hyp}$: the proof encoding is correct
  (every derivation from \texttt{hyp} has a proof tree whose
  \texttt{thFunTFor} image equals the code of the derived equation);
\item $\mathit{Consistent}\;\mathit{hyp}$: the theory cannot derive
  $O = \mathit{Pair}(O,O)$;
\item $\mathit{godelSentence}$: the G\"odel sentence, defined via diagonal
  substitution with free variables (analogous to Rose's Theorem~18).
\end{itemize}

\section{System}

\subsection{Terms and Functors}

The language has:
\begin{itemize}
\item \textbf{Terms}: $O$ (leaf), $\mathit{var}\;n$, $\mathit{ap1}\;f\;t$, $\mathit{ap2}\;g\;t_1\;t_2$
\item \textbf{Fun1} (1-ary): $I$, $\mathit{Fst}$, $\mathit{Snd}$, $\mathit{Comp}$, $\mathit{Comp2}$, $\mathit{Rec}$, $\mathit{KT}$
\item \textbf{Fun2} (2-ary): $\mathit{Pair}$, $\mathit{Const}$, $\mathit{Lift}$, $\mathit{Post}$, $\mathit{Fan}$, $\mathit{IfLf}$, $\mathit{TreeEq}$
\item \textbf{Equations}: $\mathit{eqn}\;l\;r$ (equates two terms)
\end{itemize}

Binary trees ($\mathit{Tree}$: $\mathit{lf}$ and $\mathit{nd}\;a\;b$) serve as
the data domain. $\mathit{reify} : \mathit{Tree} \to \mathit{Term}$ maps
$\mathit{lf} \mapsto O$ and $\mathit{nd}\;a\;b \mapsto \mathit{Pair}(\mathit{reify}\;a, \mathit{reify}\;b)$.

\subsection{Derivation System}

$\mathit{Deriv}\;\mathit{hyp}\;\mathit{eq}$ is an inductive type with:
\begin{itemize}
\item Computation axioms: $\mathit{axI}$, $\mathit{axFst}$, $\mathit{axSnd}$,
  $\mathit{axComp}$, $\mathit{axComp2}$, $\mathit{axLift}$, $\mathit{axPost}$,
  $\mathit{axFan}$, $\mathit{axKT}$, $\mathit{axRecLf}$, $\mathit{axRecNd}$
\item Conditional: $\mathit{axIfLfL}$, $\mathit{axIfLfN}$
\item Tree equality: $\mathit{axTreeEqLL}$, $\mathit{axTreeEqLN}$,
  $\mathit{axTreeEqNL}$, $\mathit{axTreeEqNN}$
\item Structural: $\mathit{axRefl}$, $\mathit{ruleSym}$, $\mathit{ruleTrans}$,
  $\mathit{cong1}$, $\mathit{congL}$, $\mathit{congR}$
\item Substitution: $\mathit{ruleInst}$
\item Hypothesis: $\mathit{ruleHyp}$
\item \textbf{Schema~F} (uniqueness of tree recursion):
  if $f$ and $g$ both satisfy $\mathit{Rec}\;z\;s$ (same base~$z$, same step~$s$),
  then $f = g$ on all inputs. This has 4 premises (base and step for each of $f$ and $g$).
\end{itemize}

\begin{remark}
Schema~F was corrected during this development. The original formulation
required proving $f(\mathit{Pair}(a,b)) = g(\mathit{Pair}(a,b))$ directly
(without inductive hypothesis), which is too strong. The corrected version
requires showing $f$ and $g$ satisfy the \emph{same} recursion equation---matching
Rose's original Schema~F (uniqueness of the fixed point).
\end{remark}

\section{Key Components}

\subsection{Proof Encoding (thFunTFor, Theorem~14)}

\texttt{thFunTFor} is a $\mathit{Fun1}$ (defined as $\mathit{Rec}\;O\;\mathit{thFunStep}$)
that maps proof-encoding trees to the code of the equation they prove.
The step function $\mathit{thFunStep}$ is a 26-branch dispatch handling
all derivation rules.

$\mathit{thm14E}$ (Theorem~14): if $\mathit{Deriv}\;\mathit{hyp}\;\mathit{eq}$,
then there exists a proof tree $(\mathit{pa}, \mathit{pb})$ such that
$\mathit{thFunTFor}(\mathit{Pair}(\mathit{reify}\;\mathit{pa}, \mathit{reify}\;\mathit{pb}))
= \mathit{reify}(\mathit{codeEqn}\;\mathit{eq})$.

This is proved by induction on $\mathit{Deriv}$ with one case per rule
(1402 lines in Thm14E.agda).

\subsection{Coded Substitution (substTFor, closedSubstTFor)}

\texttt{substTFor} is a $\mathit{Fun1}$ implementing coded substitution on trees,
with free variables $v_{11}$ (replacement) and $v_{12}$ (target).
$\mathit{closedSubstTFor}\;\mathit{repl}\;\mathit{tgt}$ closes these variables.

$\mathit{closedSTFCode}$ proves correctness: for any term~$l$,
$\mathit{closedSubstTFor}$ applied to $\mathit{reify}(\mathit{code}\;l)$
equals $\mathit{reify}(\mathit{codedSubst}\;\ldots)$.

\subsection{TreeEq Reflexivity via Schema~F}

$\mathit{TreeEq}(t, t) = O$ for all terms~$t$ is proved using the corrected Schema~F.
Define $f = \mathit{Comp2}\;\mathit{TreeEq}\;I\;I$ (so $f(t) = \mathit{TreeEq}(t,t)$)
and $g = \mathit{KT}\;O$ (constantly $O$). Both satisfy the recursion with base~$O$
and step $s(\mathit{orig}, \mathit{recs}) = \mathit{IfLf}(\mathit{Fst}(\mathit{recs}),
\mathit{Pair}(\mathit{Snd}(\mathit{recs}), \mathit{Pair}(O,O)))$:
\begin{itemize}
\item $f$: by $\mathit{axComp2}$, $\mathit{axTreeEqNN}$, and the recursion structure of $\mathit{TreeEq}$.
\item $g$: by $\mathit{axKT}$; the step reduces to
  $\mathit{IfLf}(O, \mathit{Pair}(O, \mathit{Pair}(O,O))) = O$ since $g$ returns $O$ on all inputs.
\end{itemize}
Schema~F then gives $f(t) = g(t)$, i.e., $\mathit{TreeEq}(t,t) = O$.

This avoids metalevel recursion on any specific tree---no heavy computation.

\section{The G\"odel Sentence and Diagonal Construction}

\subsection{Template}

The template equation uses the \emph{open} $\mathit{substTFor}$ (with free $v_{11}$, $v_{12}$):
\[
  \mathit{templateEqn} = \mathit{eqn}\;(\mathit{TreeEq}(\mathit{thFunTFor}(\mathit{var}\;0),\;
  \mathit{substTFor}(\mathit{var}\;v_1)))\;\mathit{Pair}(O,O)
\]
Using the open $\mathit{substTFor}$ keeps $\mathit{code}(\mathit{lhsT})$ small
(no nested $\mathit{Rec}$ from $\mathit{closedSubstTFor}$ or $\mathit{codeReify}$).

\subsection{Abstract Definitions}

The following are defined once in \texttt{SubstTForPrecomp.agda} (abstract, cached):
\begin{itemize}
\item $\mathit{templateCode} = \mathit{nd}\;\mathit{codeLhsT}\;\mathit{codePoo}$
  (code of the template, using abstract sub-codes)
\item $\mathit{crTC} = \mathit{code}(\mathit{reify}\;\mathit{templateCode})$
  (code of the numeral for templateCode)
\item $\mathit{godelSentence} = \mathit{substEq}\;v_1\;(\mathit{reify}\;\mathit{templateCode})\;\mathit{templateEqn}$
\item $\mathit{cGS} = \mathit{codeEqn}\;\mathit{godelSentence}$
\item $\mathit{cstf} = \mathit{closedSubstTFor}(\mathit{reify}\;\mathit{crTC})(\mathit{reify}(\mathit{natCode}\;v_1))$
\end{itemize}

\subsection{Heavy Computation (One-Time, Cached)}\label{sec:heavy}

\texttt{SubstTForPrecomp.agda} contains one step that requires \textbf{heavy computation}
(approximately 108 seconds, computed once and cached in \texttt{.agdai}):

\begin{quote}
\texttt{instForm2}: proving by \texttt{refl} that after $\mathit{ruleInst}\;v_{11}/v_{12}$
on an equation containing $\mathit{cstf}(\mathit{reify}\;\mathit{templateCode})$
and $\mathit{substTFor}(\mathit{reify}\;\mathit{templateCode})$, the
$\mathit{substTFor}$ becomes $\mathit{cstf}$. This requires Agda to traverse
$\mathit{reify}(\mathit{crTC})$ (a doubly-encoded tree with ${\sim}400$K term nodes)
with $\mathit{subst}$ as a no-op.
\end{quote}

This is the only remaining heavy computation. All other steps are lightweight
($< 8$ seconds total for all precomputation files).

\begin{remark}
The mathematical argument does not require this traversal---it follows from
the general fact that $\mathit{subst}$ is identity on closed terms. A more
refined formalization could prove this equationally (via a general ``subst
is identity on reified trees'' lemma) rather than by brute-force \texttt{refl}.
This is left for future work.
\end{remark}

\section{Proof of G\"odel II}

The proof in \texttt{GodelII.agda} (0.2 seconds, no heavy computation) follows
Rose's Theorem~18 structure:

\begin{enumerate}
\item \textbf{Assume} $\mathit{Deriv}\;\mathit{hyp}\;\mathit{godelSentence}$.

\item \textbf{thm14E} gives a proof encoding $\mathit{enc}$ with
  $\mathit{corrPf} : \mathit{thFunTFor}(\mathit{enc}) = \mathit{reify}(\mathit{cGS})$.

\item \textbf{diagFTarget} (from precomputed pieces):
  $\mathit{cstf}(\mathit{reify}\;\mathit{templateCode}) = \mathit{reify}(\mathit{cGS})$.

\item \textbf{Chain} corrPf and diagFTarget through $\mathit{reify}(\mathit{cGS})$
  as middle term: $\mathit{thFunTFor}(\mathit{enc}) = \mathit{cstf}(\mathit{reify}\;\mathit{templateCode})$.
  Note: $\mathit{reify}(\mathit{cGS})$ never appears in an equation touched by $\mathit{ruleInst}$.

\item \textbf{Instantiate} G with enc ($\mathit{ruleInst}\;0$):
  $\mathit{TreeEq}(\mathit{thFunTFor}(\mathit{enc}),\;\mathit{substTFor}(\mathit{reify}\;\mathit{TC})) = \mathit{Pair}(O,O)$.

\item \textbf{Replace} $\mathit{thFunTFor}(\mathit{enc})$ with $\mathit{cstf}(\mathit{reify}\;\mathit{TC})$
  using the chain: $\mathit{TreeEq}(\mathit{cstf}(\mathit{reify}\;\mathit{TC}),\;\mathit{substTFor}(\mathit{reify}\;\mathit{TC})) = \mathit{Pair}(O,O)$.

\item \textbf{ruleInst} $v_{11}$, $v_{12}$ (via instForm2):
  $\mathit{TreeEq}(\mathit{cstf}(\mathit{reify}\;\mathit{TC}),\;\mathit{cstf}(\mathit{reify}\;\mathit{TC})) = \mathit{Pair}(O,O)$.

\item \textbf{treeEqSelf} (via Schema~F, no heavy computation):
  $\mathit{TreeEq}(\mathit{cstf}(\mathit{reify}\;\mathit{TC}),\;\mathit{cstf}(\mathit{reify}\;\mathit{TC})) = O$.

\item \textbf{Contradiction}: $O = \mathit{Pair}(O,O)$, violating consistency.
\end{enumerate}

\section{File Structure}

\begin{center}
\begin{tabular}{lrl}
\hline
\textbf{File} & \textbf{Time} & \textbf{Role} \\
\hline
Guard/Step.agda & $<1$s & Derivation system (with corrected Schema~F) \\
Guard/Term.agda & $<1$s & Terms, functors, coding functions \\
Guard/ThFun.agda & $<1$s & Metalevel proof evaluator \\
Guard/ThFunTFor.agda & $<1$s & Object-level proof evaluator (26 branches) \\
Guard/Thm14E.agda & $<1$s & Theorem~14 (1402 lines) \\
Guard/SubstTForCorrect.agda & $<1$s & closedSubstTFor correctness \\
Guard/SubstTForPrecomp.agda & 108s & \textbf{Heavy}: all precomputed abstractions \\
Guard/SubstTForPrecomp2.agda & $<1$s & Assembly of lhsTSTF from opaque pieces \\
Guard/GodelII.agda & $<1$s & G\"odel II proof \\
\hline
\end{tabular}
\end{center}

Total: 36 files, approximately 5000 lines. The 108-second computation
is cached and not repeated.

\section{Correspondence with Rose}

\begin{center}
\begin{tabular}{ll}
\hline
\textbf{Rose} & \textbf{Ours} \\
\hline
$th(y)$ & thFunTFor \\
$se(x,y,z)$ / $st(x,y,z)$ & substTFor / closedSubstTFor \\
$nu(x)$ & codeReify (implicit via crTC) \\
$N$ & templateCode \\
Schema~F & ruleF (corrected: uniqueness of recursion) \\
Theorem~10 ($vl(x,th(y))=0$) & thm14E (ProofE) \\
Theorem~13 ($th(y) \neq e(th(z))$) & treeEqSelf + corrPf \\
Theorem~14 & thm14E \\
Theorems~15--16 & Individual cases of thm14E + congL/congR \\
Theorem~17 & diagFTarget + fixedPointEq \\
Theorem~18 & godelII \\
\hline
\end{tabular}
\end{center}

\begin{thebibliography}{1}
\bibitem{rose1961} H.\,E.\ Rose,
\emph{On the Consistency and Undecidability of Recursive Arithmetic},
Zeitschr.\ f.\ math.\ Logik und Grundlagen d.\ Math., Bd.~7, S.~124--135 (1961).
\end{thebibliography}

\end{document}
